\chapter{기호와 용어의 기준}
\label{app:notation}

\section{기호를 고정하는 이유}

기호의 일관성은 Transformer의 source와 destination 축을 뒤바꾸는 오류를 막는 첫 번째
검산 장치다. 이 책은 수식에서 column vector를 기본으로 쓰되, batch 계산의 tensor는
마지막 축을 feature로 두는 일반적인 코드 convention을 함께 표시한다.

\begin{longtable}{@{}p{0.16\textwidth}p{0.30\textwidth}p{0.43\textwidth}@{}}
\caption{핵심 기호와 차원}\label{tab:notation}\\
\toprule
기호 & 의미 & 값 또는 shape \\
\midrule
\endfirsthead
\toprule
기호 & 의미 & 값 또는 shape \\
\midrule
\endhead
$B$ & batch 크기 & 양의 정수 \\
$T$ & 현재 token sequence 길이 & 양의 정수 \\
$V$ & vocabulary 크기 & 양의 정수 \\
$L$ & Transformer block 수 & 양의 정수 \\
$d$ & residual stream 차원 & $d_{\mathrm{model}}$ \\
$d_h$ & 한 attention head의 차원 & 모델에 따라 $d/H_q$인 경우가 흔함 \\
$d_{\mathrm{ff}}$ & MLP 중간 차원 & 모델별 양의 정수 \\
$H_q$ & query head 수 & MHA에서는 보통 $H_q=H_{kv}$ \\
$H_{kv}$ & key/value head 수 & MQA에서는 1, GQA에서는 $1<H_{kv}<H_q$ \\
$t_i$ & position $i$의 token ID & $\{0,\ldots,V-1\}$ \\
$\mat{E}$ & token embedding matrix & $\R^{V\times d}$ \\
$\mat{X}^{(\ell)}$ & block $\ell$ 입구 residual stream & $\R^{B\times T\times d}$ \\
$\mat{Q},\mat{K},\mat{V}$ & query, key, value tensor & head 축을 분리하면 각각 $\shape{B,H,T,d_h}$ \\
$\mat{A}^{h}$ & head $h$의 attention weight & $\shape{B,T,T}$, 행은 destination, 열은 source \\
$\mat{W}_U$ & unembedding 또는 LM-head weight & $\R^{V\times d}$ \\
$\mat{Z}$ & vocabulary logits & $\R^{B\times T\times V}$ \\
$\theta$ & 모델의 모든 학습 parameter & vector 또는 구조화된 tensor 집합 \\
$\mathcal{L}$ & 학습 loss & scalar \\
$\Delta\vect{x}_c$ & 구성요소 $c$가 쓰는 residual update & $\R^d$ \\
$m$ & 해석 실험의 target metric & logit difference, probability, loss 등 \\
\bottomrule
\end{longtable}

\section{축과 index의 방향}

attention matrix는 destination-by-source 순서로 정의한다. 즉 $A_{i,j}$는 destination
position $i$가 source position $j$의 value를 읽는 weight다. causal decoder에서는
$j>i$인 미래 source가 mask된다.

행렬의 전치는 vector convention 때문에 달라질 수 있다. column vector를 쓰면
$\vect{q}=\mat{W}_Q^{\mathsf T}\vect{x}$로, row-vector tensor를 쓰는 PyTorch식 표기에서는
$\mat{Q}=\mat{X}\mat{W}_Q$로 적을 수 있다. 두 식의 parameter 저장 shape가 같다는
가정 없이 기계적으로 비교해서는 안 된다.

layer 번호는 문헌마다 0부터 또는 1부터 시작한다. 이 책의 수식은 embedding 뒤 residual을
$\mat{X}^{(0)}$으로 두고, block $\ell\in\{0,\ldots,L-1\}$가
$\mat{X}^{(\ell+1)}$을 만든다. 원 논문 수치를 옮길 때는 논문의 번호 convention을 그대로
표시한다.

\section{한국어와 원어 용어}

용어는 정착된 한국어가 자연스러우면 번역하고, 코드와 논문 검색에 필요한 원어는 첫
등장이나 표에서 함께 제시한다. 억지 번역으로 구현 이름을 흐리지 않기 위해 attention,
logit, embedding, probe와 checkpoint처럼 국내 기술 문헌에서 원어가 널리 쓰이는 용어는
원어를 유지한다.

\begin{longtable}{@{}p{0.24\textwidth}p{0.27\textwidth}p{0.38\textwidth}@{}}
\caption{권장 용어와 이 책에서의 의미}\label{tab:glossary}\\
\toprule
표기 & 대응 원어 & 이 책에서의 의미 \\
\midrule
\endfirsthead
\toprule
표기 & 대응 원어 & 이 책에서의 의미 \\
\midrule
\endhead
순전파 & forward pass & 고정된 입력과 parameter에서 activation과 출력을 계산하는 과정 \\
역전파 & backpropagation & chain rule로 loss의 parameter gradient를 계산하는 방법 \\
잔차 스트림 & residual stream & 각 sublayer가 update를 읽고 더하는 position별 $d$차원 상태 \\
활성값 & activation & 특정 입력에서 중간 node가 실제로 취한 수치 \\
가중치·parameter & weight, parameter & 학습되어 여러 입력에 공통으로 쓰이는 값 \\
특징·feature & feature & 의미 변수 또는 학습된 dictionary element이며 문맥별 정의가 필요함 \\
회로 & circuit & 지정한 행동을 설명한다고 가정한 계산 부분그래프 \\
개입 & intervention & activation, edge 또는 weight를 외부에서 바꾸고 결과를 재계산하는 실험 \\
복원 가능성 & decodability & 별도 readout이 내부 상태에서 변수를 예측할 수 있는 성질 \\
인과적 충실성 & causal faithfulness & 설명과 원 모델이 관련 개입에 같은 반응을 보이는 정도 \\
사실 회상 & factual recall & subject와 relation에서 attribute token을 예측하는 행동 \\
환각 & hallucination & 평가 기준상 factual claim이 근거 또는 사실과 맞지 않는 출력 \\
\bottomrule
\end{longtable}

\section{인식론적 동사의 강도}

동사의 선택은 증거 수준을 표시한다. ``계산한다''와 ``증명한다''는 정의나 수학적 귀결에,
``관찰했다''와 ``보고했다''는 특정 실험 결과에, ``시사한다''와 ``해석할 수 있다''는 대안이
남는 설명에 사용한다.

``모델이 안다'', ``생각한다''와 ``기억한다''는 조작적 정의 없이 사용하지 않는다. 예를
들어 ``모델이 Honolulu를 안다'' 대신 ``이 prompt에서 Honolulu token의 logit이 가장 높다''
또는 ``paraphrase 집합에서 정답을 생성했다''라고 쓴다. anthropomorphic 표현이 필요한
경우에는 관찰 가능한 행동을 바로 뒤에 정의한다.

\begin{studybox}
논문을 읽을 때 모든 핵심 동사에 밑줄을 긋고 ``정의'', ``관찰'', ``개입'', ``해석'' 또는
``추측''으로 분류한다. 저자의 문장도 실험 설계보다 강하면 그대로 사실 상자에 옮기지
않는다.
\end{studybox}
